%% =====================================================================
\section{Problem, Data and Metrics}
\label{sec:problem}

\paragraph{Puzzles.}
An \Sq{N} Rubik's Square is an $N\times N$ grid of cubies with a top and a bottom colour, plus side stickers on
the border. Its $2N$ moves are 180$^\circ$ flips of one row or one column; each move is its own inverse, and
flips of two rows (or two columns) commute. An \Cu{N} Rubik's Cube has $6N^2$ stickers and $3N$ slices; a
quarter, half or three-quarter turn of one slice is one move in the half-turn metric, giving $9N$ moves. Given a
reachable state, the task is to return a short move sequence that solves it. On our random benchmark a state is
\emph{solved} when every face shows one colour, in any reachable orientation (24 for the cube, 4 for the Square);
on Santa 2023 the goal is one exact state (below).

\paragraph{Clusters.}
A cubie can only ever reach the positions in its own cluster, and moves never mix clusters. Table~\ref{tab:clusters}
lists the cluster types and how many of each an $N$ has; our engine finds them automatically (\S\ref{sec:engine}).
A \Cu{20} has 2,400 stickers, 9 wing clusters and 81 center clusters, and every one must be solved.

\begin{table}[t]
\caption{Cubie cluster types. Wings combine two sticker orbits of one layer. Counts are for one
$N$; $k=\lfloor N/2\rfloor-1$.}
\label{tab:clusters}
\small
\begin{tabular}{@{}llrl@{}}
\toprule
Puzzle & Cluster type & Stickers & Count \\
\midrule
Cube & corners & 24 & 1 \\
     & middle edges & 24 & 1 if $N$ odd \\
     & wing edges & 48 & $k$ \\
     & centers ($+$, $\times$, oblique) & 24 & $((N-2)^2-[N\text{ odd}])/4$ \\
     & fixed centers & 6 & 1 if $N$ odd \\
\midrule
Square & corner & 16 & 1 \\
       & edge & 12 & $2k$ \\
       & interior & 8 & $k^2$ \\
       & edge-cross, cross, center & 6, 4, 2 & odd $N$ only \\
\bottomrule
\end{tabular}
\end{table}

\paragraph{Metrics.}
We report solution lengths after canonicalisation (merging commuting same-axis moves, \S\ref{sec:engine}) in three
metrics: the \emph{half-turn metric} (HTM), where a quarter, half or three-quarter turn of one slice counts 1; the
\emph{quarter-turn metric} (QTM), where a half turn counts 2; and the \emph{slice-turn metric} (STM), where a block
of adjacent layers turned together counts 1. On the Square all moves are flips, so HTM\,=\,QTM. The Santa 2023 score
of a cube solution is exactly its QTM: every allowed move is a quarter turn of one slice and costs 1. Random-state
results are in HTM, Santa results in QTM. To compare sizes we divide lengths by $N^2/\ln N$, the growth rate of
the diameter, or by the counting bound of \S\ref{sec:bound}.

\subsection{Data}
\label{sec:data}

\paragraph{Random states (our benchmark).}
A random state is the end of a long random canonical walk from the solved state (at least $25\times$ the number of
moves and at least 400 moves), which approximates a uniform random reachable state. Every instance is regenerated
from a seed and checked by the solvability checker, and every method sees the same instances. The main comparison
uses 10 Squares per size for $N\le12$, 5 for $13\le N\le20$, 2 cubes per size for $N\le10$ and one each at $N=15$
and $20$; the statistics of \S\ref{sec:ci} use 20 states per size for cubes $N=3$--10 and every Square size.

\paragraph{Kaggle Santa 2023 (external benchmark).}
\emph{Santa 2023 -- The Polytope Permutation Puzzle}~\cite{kaggle2023santa} was a Kaggle optimisation competition
(December 2023 to January 2024). Each of its 398 puzzles is a start state, a goal state and a list of allowed moves
given as sticker permutations; a submission lists one move sequence per puzzle, and the score is the total number of
moves, so every quarter turn counts 1. The puzzles are cubes \Cu{2} to \Cu{10}, \Cu{19} and \Cu{33}, ``wreaths'' and
``globes''. The goal is one exact state: the cube must end in the given orientation, not any rotation of it. Some
cube puzzles have other goal patterns (checkerboard faces, or every sticker its own label), and some allow a number of
wrong stickers (``wildcards''). We use the \SantaTotal{} cube puzzles whose goal has one colour per face and no
wildcards; Table~\ref{tab:santa} gives their number per size. Our adapter derives the map between Santa's sticker
order and the engine from the move permutations themselves (each Santa move must equal exactly one engine move) and
checks it by replaying KMCoders' published solutions, which all reach the exact goal in our engine. The puzzle files
are downloaded on first use and not redistributed.

\subsection{How far from optimal? A counting lower bound}
\label{sec:bound}

Every state at distance $d$ is the end point of at least one \emph{canonical} move word of length $d$, so if $c(d)$
counts such words, the radius $L$ around the goal satisfies $\sum_{d\le L}c(d)\ge|G|$, with $|G|$ the number of
reachable states. For the cube in QTM, a canonical word is a sequence of blocks; a block turns a non-empty set of
layers of one axis, each by $+1$, $-1$ or $2$ quarter turns, and consecutive blocks use different axes. A block's
cost generating function is $f(x)=(1+x)^{2N}-1$ and the words have $G(x)=1+3f/(1-2f)$, so $c(d)$ grows like $b^d$ with
$b=1/x^\ast$, $2f(x^\ast)=1$, i.e.
\begin{equation}
b=\Bigl(\bigl(\tfrac32\bigr)^{1/(2N)}-1\Bigr)^{-1},\qquad L\;\ge\;\frac{\ln|G|}{\ln b}-O(1).
\label{eq:bound}
\end{equation}
We count $|G|$ as colour states cluster by cluster: $8!\,3^7$ for the corners, $12!\,2^{11}$ for the middle edges,
$24!$ for each wing cluster and $24!/(4!)^6$ for each non-fixed center cluster (same-coloured center pieces are
indistinguishable); parity couplings change $\log_2|G|$ by about one bit per cluster and are ignored. An exact count of
words agrees with Eq.~\eqref{eq:bound} to within one move (unit tests). Because words are over-counted (different
canonical words may reach the same state), the bound stays valid. For the Square, row flips commute with row flips and
column flips with column flips, so canonical words alternate between non-empty sets of rows and of columns, $f(x)=(1+x)^N-1$
and $b=1/(2^{1/N}-1)$. Table~\ref{tab:cal} calibrates the bound where the truth is known: on the Squares whose exact
God's numbers we compute by breadth-first search, the true worst case is 1.5--1.7$\times$ the bound. Typical optimal
distances lie between the bound and the worst case. In this sense the bound turns ``how good is a solution'' into a
number for every size, including sizes where no optimal solver exists.

\begin{table}[t]
\caption{Calibrating the counting bound on the Rubik's Square: state counts and exact God's numbers (colour goal, any
of the 4 solved orientations) from breadth-first search, and the counting bound for $|G|/4$ states.}
\label{tab:cal}
\small
\begin{tabular}{@{}rrrrr@{}}
\toprule
$N$ & states & bound & God's number & ratio \\
\midrule
3 & 768 & 4 & 6 & 1.50 \\
4 & 82,944 & 6 & 10 & 1.67 \\
5 & 10,616,832 & 8 & 12 & 1.50 \\
\bottomrule

\end{tabular}
\end{table}
